
\section{Normalizar antes de comparar}
Temperatura, presión, velocidad y vibración cambian con régimen, ambiente y configuración. Comparar valores brutos puede confundir condición y degradación. Se usan correcciones y definiciones aplicables, conservando unidades y rotor de referencia. Se registran cambios de sensor, software, motor y componentes.

Un aumento de temperatura a prestación comparable puede sugerir pérdida de margen, sin identificar causa única. Se formulan hipótesis y buscan evidencias independientes. La tendencia detecta cambios; la causalidad requiere razonamiento físico.
\section{Frecuencia, orden y medida}
Orden es frecuencia dividida por giro del rotor de referencia, que se identifica en motores multieje. Seguir una orden o mantener frecuencia casi fija aporta indicios diferentes, sin diagnóstico automático. Amplitud, fase, posición, banda y condiciones también importan.

El muestreo y la ventana condicionan una señal digital; aliasing puede representar una frecuencia como otra. Se emplean equipos y procedimientos definidos, evitando reglas como «orden uno siempre significa desequilibrio». P33 verifica solo cinemática.
\section{Aceite y boroscopia como evidencias complementarias}
La muestra de aceite necesita identidad, método, punto y exposición. Concentración depende de aportación de material, reposición, filtrado y contaminación. Morfología de partículas orienta hipótesis; acciones y límites proceden de los datos pertinentes.

Boroscopia ofrece acceso y cobertura concretos. Iluminación, ángulo y referencia dimensional afectan medida. Lo no visible fuera del alcance del método no queda excluido. El libro no inventa límites de aceptación de álabe.
\section{Cerrar el ciclo de diagnóstico}
Se conserva señal, validación, hipótesis, comprobaciones, hallazgos y resultado de intervención. Si desaparece una señal, se comprueba el cambio de condiciones antes de atribuir causa. El rastro permite estudiar recurrencias y conectar daño material con decisiones documentadas.
\repaso{¿Dos señales que cambian juntas prueban una causa común?}{Pueden aumentar el interés de una hipótesis, pero también responder al régimen o a causas distintas. Se necesita evidencia discriminante.}
